\section{代码与数据说明}

本研究的数据预处理、局地化集合调整卡尔曼滤波、交叉熵航线规划器、联合嵌入预测架构时空模型、出发前预演与在线留出估值控制器、图表生成脚本、原始实验数据，以及作者与人工智能助手的研究讨论记录（已隐去涉及作者其他未发表工作的讨论内容）均开源于：\url{https://github.com/frankxucd0826-source/auv-compute-or-observe}。代码仓库中复现实验的脚本包括：用于主实验对比、开发集调参与算力相图的 \texttt{experiments/main\_eval.sh} 和 \texttt{experiments/phase.sh}，用于规划预算自评偏差实验的 \texttt{experiments/anytime\_curse.py}，用于巡测规则对比、当天卫星图真值、消融实验、1天滞后底图及单因素改动诊断的 \texttt{experiments/rev2.sh} 和 \texttt{experiments/rev2\_diag.sh}，以及用于由原始日志汇总统计置信区间与生成表格的 \texttt{experiments/main\_extra.py}、\texttt{experiments/rev2\_summary.py} 和 \texttt{experiments/rev2\_tables.py}；各脚本的具体运行命令见仓库中的 \texttt{README.md}。

实验使用的两套卫星遥感数据集均下载自美国国家海洋和大气管理局海岸观测网的公开数据服务器：
\begin{enumerate}
  \item \textbf{东海多星融合叶绿素 a 逐日数据集（2018—2026年）：}融合极轨卫星水色传感器观测并经经验正交函数缺失数据插补算法填补云区，空间分辨率为$2\text{ km}$，时间分辨率为逐日，遵循公共领域开放许可，数据集标识为 \texttt{noaacwNPPN20S3ASCIDINEOF2kmDaily}。
  \item \textbf{多尺度超分辨率海表温度逐日数据集（2018—2026年）：}由美国国家航空航天局喷气推进实验室制作（版本4.1），原始空间分辨率为$0.01^\circ$，时间分辨率为逐日，数据集标识为 \texttt{jplMURSST41}。
\end{enumerate}

\section{人工智能工具使用说明}

在本课程论文的研究与写作过程中，作者与人工智能工具的具体分工如下：
\begin{enumerate}
  \item 作者本人提出了核心研究问题，确立了研究方向与各项要求（包括第二轮修改要求：引入真实卫星遥感数据、补充基线与消融实验、采用学术期刊排版格式及提升图表质量），并对论文承担全部责任。
  \item Claude（Anthropic）完成了文献检索与核查，细化了研究选题，编写了全部代码，运行了所有数值实验，绘制了论文全部图表，并对照实验记录逐一核查了正文中的全部数值与文献引用。
  \item Gemini（Google）依据 Claude 整理的实验事实清单与写作提纲撰写了论文正文。
\end{enumerate}
